\documentclass[../../main.tex]{subfiles}

\begin{document}

\section{Riemann 积分}\label{section:ch04-Riemann积分}

我们回忆涉及 Riemann 积分的一些定义. 令 $f$ 为定义在闭有界区间 $[a,b]$ 上的有界实值函数. 令 $P=\{x_{0},x_{1},\cdots,x_{n}\}$ 为 $[a,b]$ 的一个分划, 即
\begin{align*}
a=x_{0}<x_{1}<\cdots<x_{n}=b.
\end{align*}
分别定义 $f$ 关于 $P$ 的下 Darboux 和与上 Darboux 和为
\begin{align*}
L(f,P)=\sum_{i=1}^{n}m_{i}\cdot(x_{i}-x_{i-1})
\end{align*}
与
\begin{align*}
U(f,P)=\sum_{i=1}^{n}M_{i}\cdot(x_{i}-x_{i-1}),
\end{align*}
其中, 对 $1\leqslant i\leqslant n$,
\begin{align*}
m_{i}=\inf\{f(x)\mid x_{i-1}<x<x_{i}\}\quad\text{而}\quad M_{i}=\sup\{f(x)\mid x_{i-1}<x<x_{i}\}.
\end{align*}
我们接着分别定义 $f$ 在 $[a,b]$ 上的下 Riemann 积分与上 Riemann 积分为
\begin{align*}
(R)\int_{a}^{b}f=\sup\{L(f,P)\mid P\text{ 是 }[a,b]\text{ 的分划}\}
\end{align*}
与
\begin{align*}
(R)\overline{\int}_{a}^{b}f=\inf\{U(f,P)\mid P\text{ 是 }[a,b]\text{ 的分划}\}.
\end{align*}
由于我们假设 $f$ 有界且区间 $[a,b]$ 的长度有限, 下 Riemann 积分与上 Riemann 积分是有限的. 上积分总是至少与下积分一样大, 若两个积分相等, 我们说 $f$ 在 $[a,b]$ 上是 Riemann 可积的, 并称这个公共值为 $f$ 在 $[a,b]$ 上的 Riemann 积分. 我们将它记为
\begin{align*}
(R)\int_{a}^{b}f,
\end{align*}
以暂时区别于我们在下一节考虑的 Lebesgue 积分.

定义在 $[a,b]$ 上的实值函数 $\psi$ 称为阶梯函数, 若存在 $[a,b]$ 的一个分划 $P=\{x_{0},x_{1},\cdots,x_{n}\}$ 和数 $c_{1},\cdots,c_{n}$ 使得对 $1\leqslant i\leqslant n$,
\begin{align*}
\text{若 }x_{i-1}<x<x_{i},\quad\text{则 }\psi(x)=c_{i}.
\end{align*}
观察到
\begin{align*}
L(\psi,P)=\sum_{i=1}^{n}c_{i}(x_{i}-x_{i-1})=U(\psi,P).
\end{align*}
由此以及上 Riemann 积分和下 Riemann 积分的定义, 我们推出阶梯函数 $\psi$ 是 Riemann 可积的且
\begin{align*}
(R)\int_{a}^{b}\psi=\sum_{i=1}^{n}c_{i}(x_{i}-x_{i-1}).
\end{align*}
因此, 我们可以重述下 Riemann 积分与上 Riemann 积分的定义如下:
\begin{align*}
(R)\int_{a}^{b}f=\sup\left\{(R)\int_{a}^{b}\varphi\mid\varphi\text{ 是阶梯函数且在 }[a,b]\text{ 上 }\varphi\leqslant f\right\}
\end{align*}
与
\begin{align*}
(R)\overline{\int}_{a}^{b}f=\inf\left\{(R)\int_{a}^{b}\psi\mid\psi\text{ 是阶梯函数且在 }[a,b]\text{ 上 }\psi\geqslant f\right\}.
\end{align*}

\begin{example}[Dirichlet 函数]\label{example:ch04-Dirichlet}
若 $x$ 是有理数, 设 $f(x)=1$; 若 $x$ 是无理数, 设 $f(x)=0$. 定义 $[0,1]$ 上的函数 $f$. 令 $P$ 为 $[0,1]$ 的任何分划. 根据有理数和无理数的稠密性,
\begin{align*}
L(f,P)=0\quad\text{且}\quad U(f,P)=1.
\end{align*}
因此
\begin{align*}
(R)\int_{0}^{1}f=0<1=(R)\overline{\int}_{0}^{1}f.
\end{align*}
因此 $f$ 不是 Riemann 可积的. $[0,1]$ 中的有理数集是可数的. 令 $\{q_{k}\}_{k=1}^{\infty}$ 为 $[0,1]$ 中的有理数的一个列举. 对于自然数 $n$, 在 $[0,1]$ 上定义 $f_{n}$, 若对某个 $q_{k}$ ($1\leqslant k\leqslant n$), $x=q_{k}$, 定义 $f_{n}(x)=1$, 否则, 定义 $f_{n}(x)=0$. 则每个 $f_{n}$ 是一个阶梯函数, 因此它是 Riemann 可积的. 因此, $\{f_{n}\}$ 是 $[0,1]$ 上的递增的 Riemann 可积函数序列,
\begin{align*}
\text{对所有 }n,\quad\text{在 }[0,1]\text{ 上 }|f_{n}|\leqslant 1
\end{align*}
且
\begin{align*}
\text{在 }[0,1]\text{ 上逐点地 }\{f_{n}\}\to f.
\end{align*}
然而, 极限函数 $f$ 在 $[0,1]$ 上不是 Riemann 可积的.
\end{example}

\begin{exercise}\label{exercise:ch04-01}
证明在上面 Dirichlet 函数的例子中, 在 $[0,1]$ 上 $\{f_{n}\}$ 不一致收敛于 $f$.
\end{exercise}

\begin{exercise}\label{exercise:ch04-02}
$[a,b]$ 的分划 $P'$ 称为分划 $P$ 的加细, 若 $P$ 的每个分点也是 $P'$ 的分点. 对于 $[a,b]$ 上的有界函数 $f$, 证明在加细下, 下 Darboux 和递增而上 Darboux 和递减.
\end{exercise}

\begin{exercise}\label{exercise:ch04-03}
用前一个习题证明对于闭有界区间上的有界函数, 每个下 Darboux 和不大于每个上 Darboux 和. 从这得出下 Riemann 积分不大于上 Riemann 积分.
\end{exercise}

\begin{exercise}\label{exercise:ch04-04}
假定 $[a,b]$ 上的有界函数 $f$ 在 $[a,b]$ 上是 Riemann 可积的. 证明存在 $[a,b]$ 的分划序列 $\{P_{n}\}$ 使得
\begin{align*}
\lim_{n\to\infty}[U(f,P_{n})-L(f,P_{n})]=0.
\end{align*}
\end{exercise}

\begin{exercise}\label{exercise:ch04-05}
令 $f$ 为 $[a,b]$ 上的有界函数. 假定存在 $[a,b]$ 的分划序列 $\{P_{n}\}$ 使得 $\lim_{n\to\infty}[U(f,P_{n})-L(f,P_{n})]=0$. 证明 $f$ 在 $[a,b]$ 上是 Riemann 可积的.
\end{exercise}

\begin{exercise}\label{exercise:ch04-06}
用前一个习题证明: 由于闭有界区间 $[a,b]$ 上的连续函数 $f$ 在 $[a,b]$ 上是一致连续的, 它在 $[a,b]$ 上是 Riemann 可积的.
\end{exercise}

\begin{exercise}\label{exercise:ch04-07}
令 $f$ 为 $[0,1]$ 上的递增实值函数. 对自然数 $n$, 定义 $P_{n}$ 为 $[0,1]$ 的分为 $n$ 个长度为 $1/n$ 的子区间的分划. 证明 $[U(f,P_{n})-L(f,P_{n})]\leqslant 1/n[f(1)-f(0)]$. 用习题 5 证明 $f$ 在 $[0,1]$ 上是 Riemann 可积的.
\end{exercise}

\begin{exercise}\label{exercise:ch04-08}
令 $\{f_{n}\}$ 为在闭有界区间 $[a,b]$ 上一致收敛于 $f$ 的有界函数序列. 若每个 $f_{n}$ 在 $[a,b]$ 上是 Riemann 可积的, 证明 $f$ 在 $[a,b]$ 上也是 Riemann 可积的. 即
\begin{align*}
\lim_{n\to\infty}\int_{a}^{b}f_{n}=\int_{a}^{b}f
\end{align*}
是否成立?
\end{exercise}

\end{document}